% Soluzioni dell'esercitazione TEEA 3MAT — Argomento 3: resistività, temperatura, potenza ed energia
% Per rigenerare il PDF (dalla cartella strumenti, due volte):
%   pdflatex -output-directory=../../../materiali -jobname=soluzioni-esercitazione-resistivita-potenza-energia-3MAT soluzioni.tex
% ---------------------------------------------------------------------------
%  Stile comune di esercitazione e soluzioni (pdfLaTeX).
%  Viene incluso da esercitazione.tex e soluzioni.tex con \input.
% ---------------------------------------------------------------------------
\documentclass[11pt,a4paper]{article}
\usepackage[a4paper, top=14mm, bottom=18mm, left=16mm, right=16mm, footskip=9mm]{geometry}
\usepackage[utf8]{inputenc}
\usepackage[T1]{fontenc}
\usepackage{helvet}
\renewcommand{\familydefault}{\sfdefault}
\usepackage{amsmath, amssymb, graphicx, xcolor, fancyhdr, lastpage, tabularx, booktabs, array, enumitem}
\usepackage[most]{tcolorbox}
\usepackage[siunitx, american resistors, american currents]{circuitikz}
\ctikzset{bipoles/thickness=1, current/distance=0.35}
\usetikzlibrary{calc, arrows.meta}

\definecolor{titolo}{HTML}{2A7FA8}
\definecolor{grigio}{HTML}{E9ECEF}
\definecolor{tenue}{HTML}{555555}
\setlength{\parindent}{0pt}
\setlength{\parskip}{3pt}

\newcommand{\U}[1]{\,\mathrm{#1}}
\newcommand{\virg}{{,}}
\newcommand{\riga}[1][28mm]{\underline{\hspace{#1}}}

% intestazione come nelle esercitazioni della scuola
\newcommand{\intestazione}[2]{%
  \begin{minipage}[c]{20mm}\includegraphics[width=18mm]{logo-isis.png}\end{minipage}%
  \begin{minipage}[c]{\dimexpr\linewidth-20mm}
    {\fontsize{20}{24}\selectfont\bfseries\color{titolo} #1}\par\vspace{1mm}
    {\footnotesize\bfseries Indirizzo: IP14 -- MANUTENZIONE E ASSISTENZA TECNICA \quad -- \quad Classe 3MAT}
  \end{minipage}\par\vspace{2mm}
  {\color{titolo}\rule{\linewidth}{0.8pt}}\par\vspace{3mm}
  \begin{tcolorbox}[colback=grigio, colframe=grigio, arc=0pt, left=6pt, top=4pt, bottom=4pt]
    {\Large\bfseries #2}
  \end{tcolorbox}}

\newcommand{\nota}[1]{\begin{tcolorbox}[colback=grigio!45, colframe=grigio, arc=0pt, left=6pt, top=3pt, bottom=3pt]\small #1\end{tcolorbox}}
\newcommand{\tipo}[1]{\par\vspace{2mm}{\large\color{titolo} #1}\par\vspace{-1mm}{\color{titolo!60}\rule{\linewidth}{0.4pt}}\par\vspace{1mm}}
\newcommand{\esercizio}[2]{\par\vspace{3mm}\textbf{Esercizio #1.} \textbf{#2}\par}

% alimentazione: barra con etichetta
\newcommand{\rail}[2]{\draw[line width=1.2pt] ($#1+(-0.6,0)$) -- ($#1+(0.6,0)$); \node[above] at #1 {#2};}

\pagestyle{fancy}
\fancyhf{}
\renewcommand{\headrulewidth}{0pt}
\newcommand{\piede}[1]{\fancyfoot[L]{\scriptsize I.S.I.S. Sacile-Brugnera | sede: ``Piero Della Valentina'' | #1}
  \fancyfoot[R]{\scriptsize Pagina \thepage\ di \pageref{LastPage}}}

% Schemi degli esercizi (usati sia nell'esercitazione sia nelle soluzioni)
\definecolor{rame}{HTML}{C46A2B}

\newcommand{\lineaCarico}[3]{% #1 = V, #2 = R_L, #3 = R_C
\begin{circuitikz}[scale=0.8, transform shape, line width=0.6pt]
  \rail{(0,4.4)}{$V = #1$}
  \draw (0,4.4) -- (0,4.2) to[R, l_={$R_L = #2$}, i>^={$I$}] (0,2.2) to[short, *-] (0,2.2)
        to[R, l_={$R_C = #3$}] (0,0.2) -- (0,0) node[sground]{};
  \draw (0,2.2) to[short, -o] (1.4,2.2) node[right]{X};
\end{circuitikz}}

\newcommand{\partitoreI}[3]{% #1 = I, #2 = R1, #3 = R2
\begin{circuitikz}[scale=0.8, transform shape, line width=0.6pt]
  \draw (0,4.4) node[ocirc]{} to[short, i>^={$I = #1$}] (0,3.2);
  \draw (-1.6,3.2) -- (1.6,3.2); \draw (0,3.2) node[circ]{};
  \draw (-1.6,3.2) to[short, i>^={$I_1$}] (-1.6,2.6) to[R, l_={$R_1 = #2$}] (-1.6,0.6) -- (-1.6,0.3) -- (1.6,0.3);
  \draw (1.6,3.2) to[short, i>_={$I_2$}] (1.6,2.6) to[R, l^={$R_2 = #3$}] (1.6,0.6) -- (1.6,0.3);
  \draw (0,0.3) node[circ]{} -- (0,0.1) node[sground]{};
\end{circuitikz}}

\newcommand{\sezionePista}[2]{% #1 = w, #2 = t
\begin{tikzpicture}[line width=0.6pt]
  \fill[rame!70] (0,0) rectangle (3,0.6);
  \draw[rame!60!black, line width=0.9pt] (0,0) rectangle (3,0.6);
  \draw[{Latex[length=2mm]}-{Latex[length=2mm]}] (0,-0.3) -- node[below]{$w = #1$} (3,-0.3);
  \draw[{Latex[length=2mm]}-{Latex[length=2mm]}] (3.3,0) -- node[right]{$t = #2$} (3.3,0.6);
\end{tikzpicture}}

\piede{Soluzioni esercitazione TEEA 3MAT}
\newcommand{\pas}[1]{\par\hspace*{4mm}$\displaystyle #1$\par}
\newcommand{\ris}[1]{\par\vspace{1mm}\hspace*{4mm}\colorbox{grigio}{\strut\ #1\ }\par}

\begin{document}

\intestazione{SOLUZIONI ESERCITAZIONE}{Resistività, temperatura, potenza ed energia -- Soluzioni}

\nota{I valori sono arrotondati a tre cifre significative. Nella formula $R = \rho \cdot L / S$ la lunghezza va in
metri e la sezione in mm$^2$. Un conduttore con sezione di $1\U{mm^2}$ lungo $1\U{m}$ ha un volume di $1\U{cm^3}$.}

\tipo{Resistenza di un conduttore}

\esercizio{1}{Cavo di rame: $L = 50\U{m}$, $S = 1{,}5\U{mm^2}$}
\pas{\text{a) } R_{Cu} = \rho_{Cu} \cdot \frac{L}{S} = 0{,}0175 \cdot \frac{50}{1{,}5} = 0{,}583\U{\Omega}}
\pas{\text{b) } S_{Al} = S_{Cu} \cdot \frac{\rho_{Al}}{\rho_{Cu}} = 1{,}5 \cdot 1{,}6 = 2{,}4\U{mm^2} \qquad \Big(\text{oppure } S_{Al} = \rho_{Al} \cdot \frac{L}{R} = 0{,}028 \cdot \frac{50}{0{,}583} = 2{,}4\U{mm^2}\Big)}
\pas{\text{c) } d_{Cu} = \sqrt{\frac{4 \cdot 1{,}5}{\pi}} = 1{,}38\U{mm} \qquad d_{Al} = \sqrt{\frac{4 \cdot 2{,}4}{\pi}} = 1{,}75\U{mm}}
\pas{\text{d) volume rame} = 1{,}5 \cdot 50 = 75\U{cm^3} \ \Rightarrow\ \text{peso} = 75 \cdot 8{,}96 = 672\U{g}}
\pas{\phantom{\text{d) }}\text{volume alluminio} = 2{,}4 \cdot 50 = 120\U{cm^3} \ \Rightarrow\ \text{peso} = 120 \cdot 2{,}70 = 324\U{g}}
\ris{$R = 0{,}583\U{\Omega}$ \quad $S_{Al} = 2{,}4\U{mm^2}$ \quad $d_{Cu} = 1{,}38\U{mm}$, $d_{Al} = 1{,}75\U{mm}$ \quad rame $672\U{g}$, alluminio $324\U{g}$}

\esercizio{2}{Linea aerea in alluminio: $L = 2\U{km} = 2000\U{m}$, $S = 150\U{mm^2}$}
\pas{\text{a) } R_{Al} = 0{,}028 \cdot \frac{2000}{150} = 0{,}373\U{\Omega} \qquad \text{b) } S_{Cu} = \frac{150}{1{,}6} = 93{,}8\U{mm^2}}
\pas{\text{c) } 1\U{km} \text{ di alluminio: } 150 \cdot 1000 = 150\,000\U{cm^3} \Rightarrow 150\,000 \cdot 2{,}70 = 405\,000\U{g} = 405\U{kg}}
\pas{\phantom{\text{c) }} 1\U{km} \text{ di rame: } 93{,}75 \cdot 1000 = 93\,750\U{cm^3} \Rightarrow 93\,750 \cdot 8{,}96 = 840\,000\U{g} = 840\U{kg}}
\pas{\text{d) a parità di resistenza l'alluminio pesa meno della metà e costa meno: tralicci più leggeri e distanti}}
\ris{$R = 0{,}373\U{\Omega}$ \quad $S_{Cu} = 93{,}8\U{mm^2}$ \quad alluminio $405\U{kg/km}$, rame $840\U{kg/km}$}

\esercizio{3}{Pista: $L = 80\U{mm} = 0{,}08\U{m}$, $w = 0{,}8\U{mm}$, $t = 35\U{\mu m} = 0{,}035\U{mm}$, $I = 3\U{A}$}
\pas{\text{a) } S = w \cdot t = 0{,}8 \cdot 0{,}035 = 0{,}028\U{mm^2} \qquad R = 0{,}0175 \cdot \frac{0{,}08}{0{,}028} = 0{,}05\U{\Omega}}
\pas{\text{b) } V = R \cdot I = 0{,}05 \cdot 3 = 0{,}15\U{V} \qquad P = R \cdot I^2 = 0{,}05 \cdot 3^2 = 0{,}45\U{W}}
\pas{\text{c) metà resistenza} \Rightarrow \text{sezione doppia} \Rightarrow w = 2 \cdot 0{,}8 = 1{,}6\U{mm}}
\ris{$S = 0{,}028\U{mm^2}$ \quad $R = 0{,}05\U{\Omega}$ \quad $V = 0{,}15\U{V}$ \quad $P = 0{,}45\U{W}$ \quad $w = 1{,}6\U{mm}$}

\tipo{Resistenza e temperatura, effetto Joule}

\esercizio{4}{Avvolgimento in rame: $R_0 = 4\U{\Omega}$ a $T_0 = 20\U{^\circ C}$}
\pas{\text{a) } R_{90} = R_0 \cdot [1 + \alpha \cdot (T - T_0)] = 4 \cdot (1 + 0{,}0039 \cdot 70) = 4 \cdot 1{,}273 = 5{,}09\U{\Omega}}
\pas{\text{b) } \frac{R}{R_0} = 1 + \alpha \cdot \Delta T \ \Rightarrow\ \Delta T = \frac{R / R_0 - 1}{\alpha} = \frac{4{,}8 / 4 - 1}{0{,}0039} = \frac{0{,}2}{0{,}0039} = 51{,}3\U{^\circ C} \qquad T = 20 + 51{,}3 = 71{,}3\U{^\circ C}}
\ris{$R_{90} = 5{,}09\U{\Omega}$ \quad $\Delta T = 51{,}3\U{^\circ C}$ \quad $T = 71{,}3\U{^\circ C}$}

\newpage
\esercizio{5}{Cavo di rame: $R_0 = 0{,}5\U{\Omega}$ a $20\U{^\circ C}$, $I = 16\U{A}$}
\pas{\text{a) } P_{20} = R \cdot I^2 = 0{,}5 \cdot 16^2 = 128\U{W} \qquad Q = P \cdot \Delta t = 128\U{W} \cdot 3600\U{s} = 460\,800\U{J} = 128\U{Wh}}
\pas{\text{b) } R_{60} = 0{,}5 \cdot (1 + 0{,}0039 \cdot 40) = 0{,}5 \cdot 1{,}156 = 0{,}578\U{\Omega} \qquad P_{60} = 0{,}578 \cdot 16^2 = 148\U{W}}
\pas{\text{c) } E_{20} = 128 \cdot 8 = 1024\U{Wh} \qquad E_{60} = 148 \cdot 8 = 1184\U{Wh}}
\pas{\text{la potenza aumenta: la corrente è imposta dal carico e } P = R \cdot I^2 \text{ cresce con } R}
\ris{$P_{20} = 128\U{W}$ \quad $Q = 460\,800\U{J} = 128\U{Wh}$ \quad $R_{60} = 0{,}578\U{\Omega}$ \quad $P_{60} = 148\U{W}$ \quad $E$: $1024 \to 1184\U{Wh}$}

\tipo{Computo energetico e spesa}

\esercizio{6}{Consumi dell'appartamento, $30$ giorni, $0{,}25$ euro/kWh}
\vspace{1mm}
{\renewcommand{\arraystretch}{1.3}
\begin{tabularx}{\linewidth}{@{}l c c >{\centering\arraybackslash}X >{\centering\arraybackslash}X@{}}
\toprule
\textbf{Apparecchio} & \textbf{Potenza} & \textbf{Uso al giorno} & \textbf{Energia al giorno} & \textbf{Energia al mese}\\
\midrule
bollitore & $2000\U{W}$ & $15\U{min} = 0{,}25\U{h}$ & $2000 \cdot 0{,}25 = 500\U{Wh}$ & $15\U{kWh}$\\
scaldabagno & $1200\U{W}$ & $2\U{h}$ & $1200 \cdot 2 = 2400\U{Wh}$ & $72\U{kWh}$\\
4 lampade LED & $4 \cdot 8 = 32\U{W}$ & $6\U{h}$ & $32 \cdot 6 = 192\U{Wh}$ & $5{,}76\U{kWh}$\\
computer & $200\U{W}$ & $4\U{h}$ & $200 \cdot 4 = 800\U{Wh}$ & $24\U{kWh}$\\
\midrule
\textbf{totale} & & & $\mathbf{3892\U{Wh}}$ & $\mathbf{116{,}76\U{kWh}}$\\
\bottomrule
\end{tabularx}}
\pas{\text{spesa mensile} = 116{,}76 \cdot 0{,}25 = 29{,}19 \text{ euro}}
\ris{totale $116{,}76\U{kWh}$ al mese \quad spesa $29{,}19$ euro \quad consuma di più lo scaldabagno ($72\U{kWh}$ al mese)}

\tipo{Partitori e consumi}

\esercizio{7}{$V = 24\U{V}$, $R_C = 4{,}45\U{\Omega}$, cavo di rame $1{,}5\U{mm^2}$, lunghezza totale $15 + 15 = 30\U{m}$}
\pas{\text{a) } R_L = 0{,}0175 \cdot \frac{30}{1{,}5} = 0{,}35\U{\Omega} \qquad R_L + R_C = 0{,}35 + 4{,}45 = 4{,}8\U{\Omega}}
\pas{\text{b) } V_C = V \cdot \frac{R_C}{R_L + R_C} = 24 \cdot \frac{4{,}45}{4{,}8} = 22{,}25\U{V} \qquad V_L = 24 \cdot \frac{0{,}35}{4{,}8} = 1{,}75\U{V} \qquad 22{,}25 + 1{,}75 = 24\U{V}\ \checkmark}
\pas{\text{c) } I = \frac{24}{4{,}8} = 5\U{A} \qquad P_C = 22{,}25 \cdot 5 = 111{,}25\U{W} \qquad P_L = 0{,}35 \cdot 5^2 = 8{,}75\U{W} \qquad \eta = \frac{V_C}{V} = \frac{22{,}25}{24} = 0{,}927}
\pas{\text{d) } E_C = 111{,}25 \cdot 6 = 667{,}5\U{Wh} \qquad E_L = 8{,}75 \cdot 6 = 52{,}5\U{Wh}}
\pas{\text{e) } R_L = 0{,}0175 \cdot \frac{30}{2{,}5} = 0{,}21\U{\Omega} \qquad V_C = 24 \cdot \frac{4{,}45}{4{,}66} = 22{,}9\U{V} \qquad \eta = \frac{22{,}9}{24} = 0{,}955}
\ris{$R_L = 0{,}35\U{\Omega}$ \quad $V_C = 22{,}25\U{V}$ \quad $I = 5\U{A}$ \quad $\eta = 92{,}7\,\%$ \quad $E_C = 667{,}5\U{Wh}$ \quad $E_L = 52{,}5\U{Wh}$}
\ris{con il cavo da $2{,}5\U{mm^2}$: $R_L = 0{,}21\U{\Omega}$ \quad $V_C = 22{,}9\U{V}$ \quad $\eta = 95{,}5\,\%$}

\esercizio{8}{$I = 6\U{A}$, $R_1 = 8\U{\Omega}$, $R_2 = 24\U{\Omega}$, $4\U{h}$ al giorno per $30$ giorni}
\pas{\text{a) } I_1 = I \cdot \frac{R_2}{R_1 + R_2} = 6 \cdot \frac{24}{32} = 4{,}5\U{A} \qquad I_2 = 6 \cdot \frac{8}{32} = 1{,}5\U{A} \qquad V = R_1 \cdot I_1 = 8 \cdot 4{,}5 = 36\U{V}}
\pas{\text{b) } P_1 = R_1 \cdot I_1^2 = 8 \cdot 4{,}5^2 = 162\U{W} \qquad P_2 = 24 \cdot 1{,}5^2 = 54\U{W} \qquad P_{tot} = V \cdot I = 36 \cdot 6 = 216\U{W}\ \checkmark}
\pas{\text{c) } t = 4 \cdot 30 = 120\U{h} \qquad E_1 = 162 \cdot 120 = 19\,440\U{Wh} = 19{,}44\U{kWh}}
\pas{\phantom{\text{c) }} E_2 = 54 \cdot 120 = 6480\U{Wh} = 6{,}48\U{kWh} \qquad \text{costo} = (19{,}44 + 6{,}48) \cdot 0{,}25 = 6{,}48 \text{ euro}}
\pas{\text{d) consuma di più } R_1\text{: in parallelo la tensione è la stessa e } P = V^2 / R \text{ è più grande se } R \text{ è più piccola}}
\ris{$I_1 = 4{,}5\U{A}$, $I_2 = 1{,}5\U{A}$ \quad $P_1 = 162\U{W}$, $P_2 = 54\U{W}$ \quad $E_1 = 19{,}44\U{kWh}$, $E_2 = 6{,}48\U{kWh}$ \quad $6{,}48$ euro}

\end{document}
